\section{The compiler pays only an additive cost for extra depth}
\label{sec:stable-compiler-resources}
The analog-to-digital input conversion is the only part of the comparator
whose perturbations accumulate.  Its depth is logarithmic in precision.
Once the features represent signed bits, each subsequent tanh gate
restores its output margin.  Consequently adding learner layers need not
multiply the entire comparison cost by learner depth.

\begin{lemma}[Compiler comparison resource]\label{lem:stable-compiler-resource}
Use one of the Boolean, density, or slab compilers at accuracy $\epsilon$,
with compiled width $m$, grid exponent $J$ in the continuous cases, and
$J=0$ in the Boolean case.  Write
$t=\lceil\log_2(128/\epsilon)\rceil$.
Suppose $M\ge1$ and the intended target is
$M\operatorname{clip}(g/M,[-1,1])$.
For every $N\ge m$ and every allowed learner depth $H$, there is a
sample-independent comparison law $\rho$ with
\[
 \rho\{\sup_x|z_w(x)|\le6M\}=1,\qquad
 \E_\rho\|z_w-M\operatorname{clip}(g/M,[-1,1])\|_{L^2(P_X)}^2
                 \le cM^2\epsilon,
\]
and
\begin{equation}\label{eq:stable-compiler-KL}
 \operatorname{KL}(\rho\|\mu_{N,H})\le
 c\left[(d+1)^2(J+2)^4+(d+1)(h+1)(m+1)^2
                  +(H+1)(t+2)^2\right]\Lambda^2+cM^2,
\end{equation}
where one may take
\[
 \Lambda=c(1+\alpha^{-1})\left[J+
 \log\frac{(H+2)(m+2)(d+2)(\kappa+2)(M+2)}{\epsilon}\right].
\]
For the Boolean and density cases take $\alpha=1$, and for the Boolean
case also take $\kappa=1$.  The constants are numerical except for their
dependence on the fixed scalar prior parameters.  There is no coefficient
cap on the original ReLU network.
\end{lemma}
\begin{proof}
On Boolean inputs insert one tanh layer of gain $O(\log(1/\delta))$
which restores the input signs to error at most $\delta/4$.
Its selected dense-block count is $O((d+1)^2)$.  On continuous inputs
use the quantizer with bit tolerance at most one quarter of the desired
digital tolerance.  Its prefix depth is $1+2J$, its width is at most
$d(J+2)+1$, and its total selected coefficient/vertex/layer count is at
most $c(d+1)^2(J+2)^3$.  Its coefficient logarithm is at most $c\Lambda$.

Choose the digital tolerance smaller than
$\min\{\epsilon,1/[100(m+1)]\}$, by a numerical factor where needed.
The quantizer formulas, gate count and desired accuracy ensure that
$\log(1/\delta)\le c\Lambda$.  In the density and slab cases its bad
input set has probability $O(\epsilon)$.  The ideal bit computation
approximates the clipped target in population squared error
$O(\epsilon)$, by the compiler proofs.  The fixed-precision division
step in the Boolean compiler leaves at most $2(t+1)$ final output bits.

The actual majority circuit has depth at most $b(h+1)$ and width at most
$m$.  All of its selected dense blocks, including zero coefficients,
therefore contribute at most $c(d+1)(h+1)(m+1)^2$ positions.  Padding
copies only the final $2(t+1)$ signed bits, so it contributes at most
$c(H+1)(t+2)^2$ positions.  Any unused feature has intended outgoing
coefficient zero, whose selected position is included in these counts.

Apply Lemma~\ref{lem:raw-bit-comparison} and its range-$M$ extension.
The prefix term is bounded by
$c(d+1)^2(J+2)^3[(J+2)\Lambda+\Lambda]$.
The digital term is its actual count times $c\Lambda$; the final
readout contributes $cM^2$ and logarithms already covered by $\Lambda$.
Their sum is bounded by \eqref{eq:stable-compiler-KL}.
The good-input and bad-input error bounds from the restoring-bit lemma,
combined with the compiler approximation and $|g_*|\le4$, give the stated
population error.  Its final conditioning also gives the deterministic
raw-support bound $6M$.  These are properties of a comparison law used
in the proof, not restrictions on the Bayesian prior or posterior.
\end{proof}
